\section{Trustworthiness, Limitations, and Conclusion}
\label{sec:trustworthiness}
\label{sec:limitations}

\paragraph{Trust mechanisms and human implications.}
The architecture exposes the classical counterfactual, records why a quantum branch is considered, checks every candidate with the shared classical evaluator independently of the quantum generation path, and preserves fallback. These mechanisms support auditability; they do not certify safe deployment. Customers need continuity under shocks, workers need limits on automation and review overload, and managers need accountable cost--service trade-offs. Society benefits only if adoption standards prevent performance proxies from displacing rights, welfare, and institutional legitimacy. Lower staffing cost is not automatically greater welfare; more AI or collaboration is not automatically desirable.

The stakeholder pathway is deliberately conditional. If a prospective router reduced unnecessary computation while preserving service and feasibility, organizations could redirect scarce validation and review capacity; if it failed, it could instead increase delay, opacity, and technology dependence. Workers could benefit from protected human-mandatory modes and explicit review limits, yet the same abstraction could conceal task quality, autonomy, surveillance, or unequal burden. Customers could benefit from continuity under correlated failures, but the synthetic service score is not customer welfare. These opposing pathways define what a trustworthy-AI-for-good evaluation must measure with affected stakeholders before deployment.

The same boundary sharpens the trustworthy-quantum claim. The present controls address evidence logging, candidate checking, and fallback in an idealized optimizer. They do not address malicious inputs, privacy leakage, supply-chain compromise, verifier attacks, or noise-induced QPU failures. Those are distinct security and robustness questions. Treating them as already solved would turn architectural intent into an unsupported assurance claim.

\paragraph{Scientific limits.}
The benchmark is synthetic: costs, success rates, disruptions, and penalties are not estimated from organizations or affected workers. There are no human subjects or personal data. Regional AI exposure is not fairness, collaboration is not worker benefit, and zero disallowed-mode violations is not workplace safety. Quantum evidence is limited to four retained ternary decisions, exact enumeration, noiseless statevectors, unequal depth-search budgets, no circuit decomposition, no QPU, no noise, no wall-clock or energy accounting, and no scaling test. Stronger classical full-problem and retained-core samplers could reduce or eliminate the observed signal. Bootstrap intervals describe these 32 frozen instances, not a population of workplaces.

\paragraph{Application-level path.}
Progress from L2 to L5 requires a prospective router, equal and widened classical budgets, hardware/noise studies, end-to-end time/cost accounting, operational outcomes, stakeholder-defined measures, and field-valid data with governance. Any future deployment needs human override, privacy controls, distribution-shift monitoring, and a documented accountability owner. NIST and ISO frameworks are governance references, not certifications~\citep{nist2023airmf,iso42001,iso23894}. The supported conclusion is justified, not maximal, quantum use.
